Agentic AI ecosystems increasingly rely on reusable agent skills
that combine executable code, natural-language instructions,
configuration artifacts, and external tool integrations.
These hybrid artifacts create software supply-chain risks
because malicious intent may be distributed across instructions,
code, metadata, and runtime behavior, making isolated alerts
difficult to interpret.

We evaluate three open-source scanners on a held-out test set
of 140 skills comprising 80 malicious and 60 benign samples,
including Cisco Skill Scanner with and without its
optional LLM analyzer. The default Cisco configuration achieves
a recall of 0.91 with 33 false positives among the 60 benign
skills, while the LLM-enabled configuration using
Qwen3:8B achieves a recall of 0.85 with 40 false positives.

We then present a framework that combines static findings,
safety-constrained sandbox evidence, and LLM-based reasoning
to assess malicious intent. The framework tracks the origins
of findings and relates evidence across analysis layers.
In a within-system evidence-source ablation using the same
reasoning backend, adding dynamic evidence to skill-content
reasoning increases F1 from 0.70 to 0.85. Combining static
and dynamic evidence further increases F1 to 0.89, with
five false positives among 60 benign skills. Dynamic evidence
provides the largest observed improvement.

The Full Pipeline detects 68 of 80 malicious skills,
achieving 0.88 accuracy, 0.93 precision, and 0.85 recall.
It matches the observed recall of Cisco Skill Scanner with
LLM while reducing false positives from 40 to five.
This system-level comparison does not isolate model
or architecture effects.